\section{Introduction}
\label{sec:intro}

The rise of parameter-efficient fine-tuning (PEFT), particularly Low-Rank Adaptation (LoRA) \cite{hu2021lora}, has fundamentally transformed the operational landscape of deep learning. Rather than replicating full-parameter models for every downstream task, practitioners can now train a single, frozen foundation model backbone and serve multiple specialized tasks simultaneously using lightweight, parallel adapter weights. This modular paradigm has fueled intense interest in dynamic model ensembling and serving frameworks \cite{sps2024, adamerging2023}, where the central challenge is to dynamically route each incoming test query to its corresponding task expert on-the-fly.

However, recent academic literature in this sub-area exhibits a troubling trend toward extreme over-engineering and complexity creep. To achieve accurate dynamic routing, existing state-of-the-art frameworks have progressively introduced a mountain of algorithmic and systems machinery. For instance, the Zero-Shot Centroid Alignment framework (SPS-ZCA) \cite{sps_zca2025} relies on pre-computing high-dimensional centroid statistics from a 64-sample offline calibration split, performing Unit-Norm Calibration (UNC), applying Intra-Task Dispersion Calibration (IDC) variance adjustments, tuning entropy-dependent temperatures, and fitting multi-dimensional Gaussian Mixture Models (GMMs) using Expectation-Maximization (EM) for out-of-distribution (OOD) rejection. Other recent alternatives, such as SABLE \cite{sable2025} and PFSR \cite{pfsr2025}, depend heavily on calculating cosine similarities to frozen classification-head weights. These "head-dependent" routers are fragile, as they cannot be applied in environments where classification heads are absent, non-vocabulary-aligned, or structurally mismatched.

\begin{figure}[t]
  \centering
  \includegraphics[width=\columnwidth]{results/batch_size_heterogeneity.png}
  \caption{\textbf{Resilience to Heterogeneity Collapse in Trained PyTorch Environment.} Under highly mixed heterogeneous batches, classic parametric routers (Linear Router, QWS-Merge) average ensembling coefficients over the batch to construct a single merged weight. This batch-level averaging causes their coefficients to collapse to the Uniform merging baseline ($[1/3, 1/3, 1/3]$), plummeting their accuracy to 23.96\%. In contrast, our minimalist LoRA Subspace Projection Routing (LSPR) performs sample-specific ensembling on-the-fly inside a single parallel forward pass, maintaining flat, robust accuracy (85.81\%) across all batch configurations (outperforming Uniform Merging's 23.96\% collapse).}
  \label{fig:batch_size_heterogeneity}
\end{figure}

Guided by the principle of Occam's razor, we advocate for a different path. We argue that the intrinsic mathematical geometry of the frozen PEFT weights themselves is entirely sufficient for dynamic routing. If a complex, multi-stage ensembling framework requiring offline calibration, GMM parameter fitting, and classification-head dependencies can be matched or outperformed by a clean, closed-form linear algebra projection, then the simpler method is strictly superior.

In this work, we present \textbf{LoRA Subspace Projection Routing (LSPR)}. Unlike post-hoc ensembling methods that attempt to serve arbitrary, standard pre-trained LoRA weights (which fail under LSPR due to the lack of weight-activation alignment), LSPR is a co-designed joint training-and-routing framework. LSPR pairs parameter-efficient fine-tuning with a lightweight representational autoencoding objective during training, which mathematically forces the column space of the down-projection weight matrix $A_k$ to span the task's activation subspace. At deployment time, LSPR leverages a simple geometric insight: the trained adapter matrix $A_k \in \mathbb{R}^{D \times r}$ spans an intrinsic $r$-dimensional representational subspace within the high-dimensional activation space $\mathbb{R}^D$ of the foundation model. By performing a standard, closed-form QR decomposition of the very first adapter block's down-projection matrix ($A_k = Q_k R_k$) offline, we extract an orthonormal basis $Q_k$ for each task in microseconds. At inference time, routing is accomplished by projecting early-stage activations $h_b$ onto these orthonormal bases. The L2-norm of this projection, normalized by the activation norm, represents the exact cosine of the angle between the activation vector and the task subspace. This provides a scale-invariant, highly robust alignment coordinate requiring zero hyperparameters to calibrate and zero validation data at serving time.

In the spirit of rigorous academic honesty and transparency, we evaluate LSPR within a high-fidelity, fully-trained PyTorch multi-task environment. We establish a controlled, synthetic multi-task sandbox designed as a proof-of-concept to isolate and analyze representational geometry. We train task-specific LoRA adapters on domain-shifted tasks using backpropagation, and evaluate ensembling and serving performance on the physically learned weights and activations.

Under this challenging setting, LSPR naturally resolves several critical flaws of prior SOTA methods:
\begin{enumerate}
    \item \textbf{The Early-Layer Routing Paradox:} Head-dependent methods (SABLE, PFSR) route using classification-head outputs. Consequently, they must execute all task-specific adapters throughout the model's entire depth before routing can occur, creating representational conflicts in early layers. LSPR routes at an early stage (such as Layer 3 of the backbone) using activation-space projections. This allows early layers of the backbone to be executed task-agnostically with no adapters, preserving expert capacity in subsequent layers. (Note that in our proof-of-concept Isolating Coordinate Sandbox simulation, the physical backbone models a single-layer projection for clarity and representation-space isolation; the multi-layer, early-layer routing advantages represent a structural property of LSPR when deployed on deep multi-layer Transformer backbones.)
    \item \textbf{Immunity to Heterogeneity Collapse:} Under mixed heterogeneous batches, classic parametric routers (e.g., Linear Router, QWS-Merge) experience "heterogeneity collapse" where ensembling coefficients average out to a uniform blend, causing performance to drop catastrophically (as shown in Figure~\ref{fig:batch_size_heterogeneity}). LSPR executes a vectorized, single-pass parallel forward pass that applies sample-specific activation blending on-the-fly, completely bypassing batch-level averaging.
    \item \textbf{Flat Latency Scaling:} While sequential micro-batching solutions (such as MBH \cite{pfsr2025}) scale linearly in latency due to multiple sequential DRAM weight-reloads, LSPR serves mixed batches in a single parallel pass, achieving outstanding physical latency scaling on edge CPUs.
\end{enumerate}

Our core contributions can be summarized as follows:
\begin{itemize}
    \item \textbf{Mathematical Simplicity:} We introduce LSPR, which replaces the highly over-engineered pipelines of prior SOTA (calibration datasets, GMMs, UNC, classification-head dependencies) with a single, elegant QR decomposition and projection.
    \item \textbf{Robust Domain-Shifted Performance:} Under both homogeneous and heterogeneous streams with non-trivial domain shifts within our controlled synthetic proof-of-concept, LSPR recovers 85.81\% Joint Mean accuracy (exceeding Uniform Merging by 61.85\% and perfectly recovering the expert ceiling and SPS-ZCA SOTA) with zero trainable parameters and zero task-specific calibration data at deployment time, enabled by our joint training-and-routing paradigm.
    \item \textbf{Head-Free OOD Rejection:} LSPR's zero-shot projection energy detects out-of-distribution (OOD) tasks with an outstanding AUROC of 0.9755 under domain shifts, far surpassing SABLE and matching EM-fitted GMMs without fitting any parametric density model.
    \item \textbf{Computational Efficiency:} We show that LSPR runs inside a single-pass vectorized execution layout, eliminating DRAM weight-loading overhead to deliver flat, highly efficient physical execution scaling on host CPUs.
\end{itemize}
